\section{A finite rational descent test for integral-slope germs}

This section makes the geometric multiplicity criterion CC-AG09
effective in a specified sector.  It is a deduction from the
checked Milne classification, elementary ideal theory, and
Kronecker's root-of-unity criterion.  The category remains the
abelian, Artin, and Tate numerical category.  No novelty claim
or extension to arbitrary motives is made.

\subsection{CC-AG11: a cyclic multiplicity monoid with a finite root test}

Let \(\xi\in\mathcal W_p\) have weight \(w\in\mathbb Z\).
Choose a stabilized representative
\[
          \xi=[\pi,b],\qquad \pi\in K,\quad b\ge1,\qquad
          K=K_\xi=\mathbb Q(\pi^N)\quad(N\ge1).
                                                               \tag{CC-AG.50}
\]
Such a representative exists by (CC-AG.24).
For each place \(v\mid p\) of \(K\), normalize its
valuation to have image \(\mathbb Z\), put
\(e_v=\operatorname{ord}_v(p)\), and define
\[
       \lambda_v=\frac{\operatorname{ord}_v(\pi)}{b e_v}.
                                                               \tag{CC-AG.51}
\]
Assume throughout this theorem that every
\(\lambda_v\) is an integer.  Define the fractional ideal
\[
                   J=\prod_{v\mid p}\mathfrak p_v^{e_v\lambda_v}
                             \quad\text{of }K.                \tag{CC-AG.52}
\]
Negative exponents are allowed.  In particular,
\((\pi)=J^b\).

Let
\[
 \mathscr E_\xi=
 \left\{n\in\mathbb Z:\ \text{there is }\gamma\in K^\times
       \text{ with }(\gamma)=J^n
       \text{ and }|\sigma\gamma|=p^{nw/2}
                    \text{ for every }\sigma:K\hookrightarrow\mathbb C
 \right\}.
                                                               \tag{CC-AG.53}
\]
Then \(\mathscr E_\xi=h\mathbb Z\) for a positive integer
\(h\) dividing \(b\).
This integer can be found by the finite test
\[
 h=\min\left\{n>0:n\mid b,\quad
       X^b-\zeta\pi^n\text{ has a root in }K
                      \text{ for some }\zeta\in\mu(K)\right\}.
                                                               \tag{CC-AG.54}
\]
Here \(\mu(K)\) is the finite group of roots of unity in \(K\).
Thus only divisors of a given exponent and roots of finitely
many specified polynomials over a fixed number field are
required.

Let \(S_\xi\) denote the rational simple object of the
\(\Gamma\)-orbit of \(\xi\) over
\(\overline{\mathbb F}_p\).  For every \(a\ge1\), its
exact multiplicity monoid for descent to \(\mathbb F_{p^a}\)
is
\[
 H_{\overline{\mathbb F}_p}(O_\xi)
                 =\frac{h}{\gcd(h,a)}\,\mathbb N.
                                                               \tag{CC-AG.55}
\]
Equivalently, for every \(r\ge0\),
\[
 S_\xi^{\oplus r}\text{ is a base extension from }
             \mathcal A_{p^a}
       \quad\Longleftrightarrow\quad
                  \frac{h}{\gcd(h,a)}\mid r.                 \tag{CC-AG.56}
\]
Every relevant target and source simple has Schur index one.
One source simple already realizes the least positive
multiplicity.  More precisely, choose \(\gamma\) in
(CC-AG.53) for \(n=h\), and put \(g=\gcd(h,a)\).
Any algebraic root
\[
                     \alpha^{h/g}=\gamma^{a/g}               \tag{CC-AG.57}
\]
is a Weil \(p^a\)-number of weight \(w\), has germ \(\xi\),
and satisfies
\[
                [\mathbb Q(\alpha):K]=h/g.                  \tag{CC-AG.58}
\]
This gives an actual finite-degree construction of a
minimal preimage.

\paragraph{Complete proof.}
The set in (CC-AG.53) contains zero and is closed under
addition and negation, by multiplying and inverting its
witnesses.  The element \(b\) belongs to it, witnessed
by \(\pi\), since (CC-AG.51)--(CC-AG.52) give
\((\pi)=J^b\) and \(\pi\) has the stated Weil absolute
values.  Hence \(\mathscr E_\xi\) is a nonzero subgroup
of \(\mathbb Z\), equal to \(h\mathbb Z\) with \(h\ge1\)
and \(h\mid b\).

We verify the finite test in both directions.
If \(\gamma\) witnesses \(n\in\mathscr E_\xi\), then
\[
                              u=\gamma^b/\pi^n
                                                               \tag{CC-AG.59}
\]
has trivial fractional ideal, and every complex conjugate
has absolute value one.  Thus \(u\) is an algebraic
integer and a unit.  Kronecker's criterion makes it
a root of unity, and it lies in \(K\), so
\(u\in\mu(K)\).  Conversely, suppose
\(\gamma\in K\) satisfies
\(\gamma^b=\zeta\pi^n\) for \(\zeta\in\mu(K)\).
It is nonzero.  Taking fractional ideals and using
their torsion-free group gives
\[
                    (\gamma)^b=J^{bn}
                       \quad\Longrightarrow\quad
                    (\gamma)=J^n.
 \]
Taking the absolute value at every complex embedding
gives \(|\sigma\gamma|=p^{nw/2}\).
Therefore \(\gamma\) witnesses membership in
\(\mathscr E_\xi\).  Since its positive generator is a
divisor of \(b\), testing the positive divisors of \(b\)
is sufficient and proves (CC-AG.54).

The precise root-of-unity input is Milne (2020),
Proposition~5.5 and Corollary~5.6, printed p.~87.
The elementary reason is that the minimal polynomials
of the powers of an algebraic integer whose conjugates
have absolute value one have uniformly bounded degree
and integer coefficients bounded by elementary symmetric
functions of numbers of modulus one.  Only finitely
many polynomials, and hence finitely many powers, can
occur.  Two powers coincide, which proves torsion.
The same bounded-polynomial argument proves the
finiteness of \(\mu(K)\), and supplies an explicit finite
way to determine it if it is not already known.  Put
\[
 R_D=\sum_{e=1}^{D}e\prod_{j=1}^{e}
                    \left(2\binom ej+1\right),
 \qquad D=[K:\mathbb Q],\qquad
 M_D=\operatorname{lcm}(1,2,\ldots,R_D).
\]
There are at most \(R_D\) complex numbers whose monic
minimal polynomials have degree at most \(D\) and the
coefficient bounds just described.  All the distinct
powers of a root of unity in \(K\) are among them, so
its order is at most \(R_D\).  Consequently \(\mu(K)\)
is exactly the set of roots in \(K\) of
\(X^{M_D}-1\).  This deliberately large bound requires
no search without a stopping condition.
Formula (CC-AG.54) is therefore a finite collection
of exact number-field root questions.

We next justify the Schur-index assertion, including
the archimedean places.  If \(K\) has a real embedding,
the Weil condition on \(\pi\) implies
\(\pi^2=p^{bw}\): its image under that real embedding
has square \(p^{bw}\), and the embedding is injective.
The stabilized condition (CC-AG.50) then gives
\(K=\mathbb Q(\pi^2)=\mathbb Q\).
Its single \(p\)-adic slope is \(w/2\), so the
integral-slope assumption forces \(w\) to be even.
Otherwise \(K\) has no real embeddings.
Consequently every real-place invariant of \(S_\xi\)
vanishes.  Its \(p\)-adic invariant at \(v\) is
\[
                       -\lambda_v[K_v:\mathbb Q_p],
\]
an integer, and all other invariants vanish by
(CC-AG.25).  Hence \(d_\xi=1\).

Take any source label \(\alpha\in W(p^a)\) mapping to
the orbit of \(\xi\).  Conjugating it if necessary,
assume that \([\alpha,a]=[\pi,b]\).  Some \(N\ge1\)
then satisfies
\[
                           \alpha^{bN}=\pi^{aN}.             \tag{CC-AG.60}
\]
Set \(L=\mathbb Q(\alpha)\) and \(d=[L:K]\).
The stabilized condition gives
\(K=\mathbb Q(\pi^{aN})\subset L\), so this degree
is well-defined.
Taking ideals in \(L\) in (CC-AG.60) and using
torsion-freeness gives
\[
                              (\alpha)=J^a\mathcal O_L.
                                                               \tag{CC-AG.61}
\]
At a place \(z\) of \(L\) above a place \(v\mid p\)
of \(K\), the corresponding normalized slope of
\(\alpha\), relative to \(p^a\), is exactly
\(\lambda_v\); this follows either from
(CC-AG.60) or from (CC-AG.61), accounting for the
ramification index in \(L/K\).
It is an integer.  If \(L\) has a real embedding,
its restriction makes \(K\) real, so \(w\) is even
by the preceding paragraph.
Thus the local-invariant formula also gives
\(d_\alpha=1\).

Now apply the relative field norm to \(\alpha\).
The ideal norm identities, which hold without a
Galois assumption, give
\[
                  \bigl(N_{L/K}(\alpha)\bigr)=J^{ad}.
                                                               \tag{CC-AG.62}
\]
For each embedding \(\sigma:K\hookrightarrow\mathbb C\),
the conjugate of this norm is the product over the
\(d\) embeddings of \(L\) extending \(\sigma\).
Each factor has absolute value \(p^{aw/2}\).
Therefore the norm is a witness for
\(ad\in\mathscr E_\xi=h\mathbb Z\).
It follows that
\[
                           h/\gcd(h,a)\mid d.                 \tag{CC-AG.63}
\]
The ideal and field norm formulas used here are Milne
(2020), Proposition~4.1(a),(c), p.~69, and
Proposition~2.19, p.~31, respectively.
Because both indices are one, (CC-AG.35)--(CC-AG.41)
identify the transition multiplicity of this source
simple with \(d\).
Thus every allowed multiplicity is divisible by
\(h/\gcd(h,a)\).

It remains to construct that least positive
multiplicity.  Choose \(\gamma\) witnessing exponent
\(h\) and let \(\alpha\) satisfy (CC-AG.57).
Every conjugate of \(\alpha\) has absolute value
\(p^{aw/2}\).  Valuations away from \(p\) vanish,
because those of \(\gamma\) vanish, so
\(\alpha\in W(p^a)\).
Equation (CC-AG.59) for \(n=h\) shows that
\([\gamma,h]=[\pi,b]=\xi\).
Moreover (CC-AG.57), raised to the \(g\)-th power,
is \(\alpha^h=\gamma^a\), so
\([\alpha,a]=[\gamma,h]=\xi\).

The germ centre is \(K\).  Since \(\gamma\in K\),
every positive power of \(\gamma\) generates \(K\):
its field contains the eventual stabilized germ
centre and is contained in \(K\).
Consequently
\(\gamma^{a/g}=\alpha^{h/g}\) proves
\(K\subset\mathbb Q(\alpha)\).
The defining equation bounds this extension degree
above by \(h/g\), whereas (CC-AG.63) says it is a
positive multiple of \(h/g\).
It must therefore equal \(h/g\), proving
(CC-AG.58).  Its source simple realizes exactly
\(h/g\) copies of \(S_\xi\).
Direct sums realize all nonnegative multiples, and
the preceding divisibility excludes every other
multiplicity.  Equations (CC-AG.55)--(CC-AG.56)
follow. \(\square\)

\paragraph{Ordinary imaginary-quadratic special case.}
Suppose \(K\) is imaginary quadratic, \(w=1\), and the
germ is ordinary, so that \(p\) splits and the two
slopes are zero and one.  Let \(\mathfrak p\) be the
prime where the slope is one.
Then \(J=\mathfrak p\), and
\[
                           h=\operatorname{ord}([\mathfrak p])
                                  \quad\text{in }\operatorname{Cl}(K).
                                                               \tag{CC-AG.64}
\]
Indeed, a Weil generator of \(J^n\) makes \(J^n\)
principal.  Conversely, if
\(J^n=(\gamma)\), then
the positive field norm of \(\gamma\) is \(p^n\);
the two complex conjugates therefore each have
absolute value \(p^{n/2}\).  Thus this generator
satisfies (CC-AG.53).  The two exponent subgroups
are identical, proving (CC-AG.64).
The exact multiplicity for descent to \(\mathbb F_{p^a}\)
is consequently
\[
             \frac{\operatorname{ord}([\mathfrak p])}
                  {\gcd(\operatorname{ord}([\mathfrak p]),a)}.
\]
For the germ of CC-AG10, this order is two and
\(a=1\), recovering its exact even-multiplicity
obstruction.  For \(a=2\) the least multiplicity
is one.

\paragraph{Exact boundary.}
The integral-slope hypothesis makes every Schur index
in this proof one.  With fractional slopes the local
indices can vary with the source centre field, and
the transition multiplicity is the full expression
\([K_\alpha:K_\xi]d_\alpha/d_\xi\).
Neither the cyclic-monoid formula (CC-AG.55) nor the
root-degree construction is asserted in that remaining
case.  This result removes the previously unbounded
root-order search precisely in the stated sector.
It does not provide new cross-correspondences for an
arbitrary variety, so the separate arbitrary-motive
Tate gate is unchanged.

\subsection{Primary references}

\begin{itemize}
\item Milne, J.~S. (1994). Motives over finite fields. In
U.~Jannsen, S.~Kleiman, \& J.-P.~Serre (Eds.),
\emph{Motives} (Proceedings of Symposia in Pure Mathematics,
Vol.~55, Part~1, pp.~401--459). American Mathematical
Society. Corrected author version dated November~2, 2019:
\url{https://www.jmilne.org/math/xnotes/Seattle1.pdf}.
Exact classification and local-index locators are
Theorems~2.15--2.16, p.~18, and Theorems~2.18--2.19,
p.~20, with the geometric-Frobenius normalization
on p.~17.
\item Milne, J.~S. (2020, July~19).
\emph{Algebraic number theory} (Version~3.08)
[Author course notes].
\url{https://www.jmilne.org/math/CourseNotes/ANT.pdf}.
Exact inputs: Proposition~2.19, p.~31;
Theorem~3.20, p.~53; Proposition~4.1(a),(c),
p.~69; Proposition~5.5 and Corollary~5.6, p.~87.
\end{itemize}
